\documentclass[11pt,a4paper]{article}
\usepackage[margin=27mm]{geometry}
\usepackage{amsmath,amssymb}
\usepackage[T1]{fontenc}
\setlength{\emergencystretch}{2em}
\title{A dilation obstruction to the displayed vertex-volume bound\\
\large Disproof of conjecture 00000007801}
\author{ziangni-sys}
\date{4 October 2026}
\begin{document}
\maketitle
The source defines $G_k(K)$ as the mean volume of the
$k$-simplices spanned by the vertices of a $d$-dimensional convex
body, and explicitly asserts
\[
 G_k(K)\leq\binom dk V(S),
\]
where $S$ is the standard simplex.
Both language versions give this displayed inequality.
In dimension $d=k=1$, an actual interval of length two disproves it.
There is no volume-normalization hypothesis on $K$ in the source.
This concerns the stated formula, rather than a separately
formulated classical Godbersen inequality.

\textbf{Actual full-dimensional convex body.}
Use the ordinary real Euclidean line and set
\[
 K=[0,2],\qquad S=[0,1].
\]
These are compact convex bodies. The interval $K$ has nonempty
interior: the point $1$ belongs to $(0,2)$.
Its actual ambient real dimension is one.
The standard one-simplex $S$ has ordinary Lebesgue volume one.

\textbf{Exhaustive vertices of $K$.}
Its extreme points are exactly $0$ and $2$.
To verify the left endpoint, suppose
\[
 0=ax+by,\qquad a,b>0,\quad a+b=1,\quad x,y\in[0,2].
\]
Both terms are nonnegative, so $ax=0$ and therefore $x=0$;
likewise $y=0$. Thus $0$ is extreme.
For the right endpoint, an analogous representation
$2=ax+by$ gives
\[
 a(2-x)+b(2-y)=0.
\]
Both terms are nonnegative, so $x=y=2$.

Conversely, if $0<z<2$, then
\[
 z=(1-z/2)\,0+(z/2)\,2
\]
is a strict convex combination of distinct endpoints.
Hence $z$ is not extreme.
The complete vertex set is therefore $\{0,2\}$.

\newpage
\textbf{Complete mean of actual vertex simplices.}
A one-simplex requires two distinct vertices. There is exactly
one unordered two-element subset of $\{0,2\}$, namely the
whole vertex set.
The corresponding simplex is its genuine convex hull:
\[
 \operatorname{conv}\{0,2\}=[0,2]=K.
\]
Therefore the mean is a sum over one actual simplex divided by
the number of such simplices:
\[
 G_1(K)=\frac{V(\operatorname{conv}\{0,2\})}{1}
       =V([0,2])=2.
\]
Using ordered distinct vertex pairs would give two identical
volumes and the same mean, so no averaging convention affects
this example.

\textbf{Failure of the displayed inequality.}
The source's right-hand side is
\[
 \binom11V(S)=1\cdot V([0,1])=1.
\]
Thus its claimed inequality is the false statement
\[
 2\leq1.
\]
The source also displays the proposed coefficient
\[
 c_{d,k}=\frac{d!}{(d-k)!\,k!}\frac{V(S)}{V(K)}.
\]
Here this coefficient equals $1/2$, so the original form
$G_1(K)\leq c_{1,1}V(K)$ again gives $2\leq1$.

\textbf{Scope and scaling.}
Dilation makes the defect transparent: for $K=[0,L]$ the sole
vertex one-simplex has volume $L$, while the fixed standard
simplex continues to have volume one.
The actual counterexample uses $L=2$ and satisfies all stated
convex-body conditions.
The special middle-dimensional stability assertion is a
separate conjunct and is not needed.
No disproof of a different normalized or mixed-volume
formulation is claimed.

\textbf{Formal correspondence and reproduction.}
Lean uses the actual real interval, compactness, convexity,
interior and real dimension. It proves the complete set of
Mathlib extreme points, a bijective finite enumeration of the
vertices, and exhaustive enumeration of every two-vertex subset.
Each simplex is the actual convex hull of its selected vertices.
Its volume is actual Lebesgue measure converted to a real number.
The finite mean, both actual interval volumes, the displayed
coefficient and the failure of both displayed inequalities are
proved.

Run \texttt{lake build} in \texttt{lean/} with Lean 4.19.0 and
public Mathlib revision:
\begin{center}\small
\texttt{c44e0c8ee63ca166450922a373c7409c5d26b00b}
\end{center}
The project prints the relevant axiom audits.
\end{document}
